% ============ §9 复习课：连词、代词格与考试题型 ============
\hdr{§9}{复习课：连词、代词格与考试题型}

\callout{\textbf{这一节的坐标}：ПОВТ p2（分值表）、p3（Как зову́т 与人称代词的格）、p31--p32（连词）、p33（组句）。\\
\textbf{这一节要解决的事}：把最后一份复习课件（Урок 9 ПОВТОРЕНИЕ）里、前面八节没讲过的两件语法补上，再把它的练习整理成一张「考试会怎么考」的清单。}

\subhdr{9.1　连词：и／а／но 说的是三种关系}

课件 ПОВТ p31 把连词单独列成一页，p32 给了三道填空。三个连词的差别不在词义，而在\textbf{前后两句是什么关系}：

\begin{minipage}{\linewidth}
\tabtop
\begin{tabular}{@{}>{\raggedright\arraybackslash}p{0.13\linewidth}>{\raggedright\arraybackslash}p{0.20\linewidth}>{\raggedright\arraybackslash}p{0.57\linewidth}@{}}
\thead{连词} & \thead{关系} & \thead{课件原例（ПОВТ p31）} \\
\end{tabular}
\tabmid
\begin{tabular}{@{}>{\raggedright\arraybackslash}p{0.13\linewidth}>{\raggedright\arraybackslash}p{0.20\linewidth}>{\raggedright\arraybackslash}p{0.57\linewidth}@{}}
и & 并列，两件事都成立 & Э́то моя́ подру́га, \textbf{и} э́то то́же моя́ подру́га. \\
а & 对比，两件事不一样 & Э́то моя́ но́вая ру́чка, \textbf{а} э́то твоя́ ста́рая ру́чка. \\
но & 转折，后一件与预期相反 & Э́та пло́щадь о́чень ста́рая, \textbf{но} о́чень краси́вая. \\
\end{tabular}
\tabbot
\end{minipage}

\textbf{这张表在说什么}：\textbf{и} 和 \textbf{а} 都可能被译成「和／而」，区别只在\textbf{前后是不是同一个东西}：同向叠加用 и，一正一反摆在一起用 а。\textbf{но} 更强：它后面必须是一个「本来不该这样、但确实这样」的事实。

\textbf{逐行读法}：第一行前后都是「我的好朋友」，所以用 и；第二行前面是新的、后面是旧的，两枝笔对着摆，所以用 а；第三行「很旧」本来会让人以为「不好看」，结果「很好看」，用 но 把这个预期掰回来。

\warnbox{\textbf{边界}：\textbf{а} 不是「但是」。真要表达「虽然……但是……」，用 \textbf{но}；用 а 会把句子变成「对比两件事」，意思跑掉。}

\subhdr{9.2　另外两个连接词：потому́ что 与 что}

ПОВТ p31 的同一页还给了两组「带从句」的连词，它们和 и／а／но 不是一个层级——后面跟的是\textbf{一个完整的句子}：

\begin{minipage}{\linewidth}
\tabtop
\begin{tabular}{@{}>{\raggedright\arraybackslash}p{0.20\linewidth}>{\raggedright\arraybackslash}p{0.18\linewidth}>{\raggedright\arraybackslash}p{0.52\linewidth}@{}}
\thead{连接词} & \thead{回答什么} & \thead{课件原例} \\
\end{tabular}
\tabmid
\begin{tabular}{@{}>{\raggedright\arraybackslash}p{0.20\linewidth}>{\raggedright\arraybackslash}p{0.18\linewidth}>{\raggedright\arraybackslash}p{0.52\linewidth}@{}}
потому́ что & Почему́?（为什么） & Мне нра́вится Пеки́н, \textbf{потому́ что} э́то большо́й совреме́нный го́род. \\
что & Что?（什么） & Я ду́маю, \textbf{что} Пеки́н о́чень краси́вый го́род. \\
\end{tabular}
\tabbot
\end{minipage}

\textbf{这张表在说什么}：两个词都译「因为／……」之外还要看\textbf{主句动词}。\textbf{потому́ что} 只能接在「喜欢、高兴、迟到」这类有原因可讲的主句后面；\textbf{что} 接在「说、想、知道、觉得」这类\textbf{引述}动词后面（говори́ть／ду́мать／знать／понима́ть）。

\textbf{逐行读法}：第一句主句是「我喜欢北京」，后面解释原因，用 потому́ что；第二句主句是「我认为」，后面是「认为的内容」，用 что。做题时先看主句动词，再看该用哪个。

\aside{课件的 p32 把这两组做成了选择题：前三题在 потому́ что／что／почему́ 里选，后三题在 и／а／но 里选。选择题的做法就是先判「主句动词」或「前后关系」，再选词——不要凭语感。}

\subhdr{9.3　人称代词的格：Меня́ зову́т 里的那个 меня́ 是哪来的}

ПОВТ p3 把一个表摆在 Как зову́т 旁边，左列是第一格（主语），右列是\textbf{第四格}（宾语）：

\begin{minipage}{\linewidth}
\tabtop
\begin{tabular}{@{}>{\raggedright\arraybackslash}p{0.16\linewidth}>{\raggedright\arraybackslash}p{0.16\linewidth}>{\raggedright\arraybackslash}p{0.16\linewidth}>{\raggedright\arraybackslash}p{0.16\linewidth}>{\raggedright\arraybackslash}p{0.16\linewidth}@{}}
\thead{第一格} & \thead{第四格} & \thead{第一格} & \thead{第四格} & \thead{备注} \\
\end{tabular}
\tabmid
\begin{tabular}{@{}>{\raggedright\arraybackslash}p{0.16\linewidth}>{\raggedright\arraybackslash}p{0.16\linewidth}>{\raggedright\arraybackslash}p{0.16\linewidth}>{\raggedright\arraybackslash}p{0.16\linewidth}>{\raggedright\arraybackslash}p{0.16\linewidth}@{}}
я & меня́ & мы & нас & 「我／我们」被问、被叫 \\
ты & тебя́ & вы & вас & 「你／您」被问、被叫 \\
он & его́ & она́ & её & 与物主代词同形 \\
они́ & их & — & — & 与物主代词同形 \\
\end{tabular}
\tabbot
\end{minipage}

\textbf{这张表在说什么}：\textbf{зову́т} 后面要的是「被叫的那个人」，所以代词不能停在第一格。中文说「我叫安娜」，「我」是主语；俄语说 \textbf{Меня́ зову́т А́нна}，字面是「（他们）叫我安娜」，「我」变成了宾语，于是用 меня́。

\textbf{逐行读法}：问句 \textbf{Как тебя́ зову́т?}／\textbf{Как Вас зову́т?} 用的是同一个第四格；回答 \textbf{Меня́ зову́т …}。说第三个人时 \textbf{Его́ зову́т Анто́н}／\textbf{Её зову́т Ма́ша}——注意 его́、её、их 这三个形式既是第四格，也是物主代词「他的／她的／他们的」，\S3.3 已经用过一次。

\warnbox{\textbf{边界}：不要写 *Я зову́т А́нна。这是本课最高频的错误——错在把「我」留在第一格。}

\subhdr{9.4　考试会怎么考：复习课件给的题型清单}

ПОВТ p2 是评分表：\textbf{Ⅰ 60 分 ＋ Ⅱ 25 分 ＋ Ⅲ 15 分 ＝ 100 分}。也就是说考试分三部分；复习课件后面每一页的练习，就是这三部分会用到的题型。按页面顺序整理如下：

\begin{minipage}{\linewidth}
\tabtop
\begin{tabular}{@{}>{\raggedright\arraybackslash}p{0.05\linewidth}>{\raggedright\arraybackslash}p{0.24\linewidth}>{\raggedright\arraybackslash}p{0.13\linewidth}>{\raggedright\arraybackslash}p{0.48\linewidth}@{}}
\thead{\#} & \thead{题型} & \thead{课件页} & \thead{怎么答（关键一步）} \\
\end{tabular}
\tabmid
\begin{tabular}{@{}>{\raggedright\arraybackslash}p{0.05\linewidth}>{\raggedright\arraybackslash}p{0.24\linewidth}>{\raggedright\arraybackslash}p{0.13\linewidth}>{\raggedright\arraybackslash}p{0.48\linewidth}@{}}
1 & 疑问词选择：Что／Кто／Где／Как／Когда́ & p9 & 先看答句是「人／东西／地点／方式／时间」，再选词 \\
2 & 连词填空：потому́ что／что／почему́ & p32 & 看主句动词是「喜欢」还是「认为」 \\
3 & 连词填空：и／а／но & p32 & 看前后是「同一件事」「两件对照」还是「转折」 \\
4 & 名字与格：Меня́／Тебя́／Его́ зову́т & p3 & 主语位置换成第四格，зову́т 不变 \\
5 & 名词的性与数（кто／что、单复数） & p5--p12 & 先判性，再按硬／软词干加词尾 \\
6 & 指示代词＋形容词一致（э́тот／са́мый） & p13--p16, p19 & 三处词尾同进同退，看名词的性数 \\
7 & нра́виться 的单复数选择 & p20--p21 & 看主语是单数还是复数：нра́вится／нра́вятся \\
8 & 动词变位与过去时填空 & p22--p30 & 先定变位组，再按人称／性数填 \\
9 & 问答配对（Как дела́? 一类） & p4 & 问句定主题，答句只用课上给的那几个词 \\
10 & 组句（词序） & p33 & 先找主语，再找动词，副词放动词前 \\
\end{tabular}
\tabbot
\end{minipage}

\textbf{这张表在说什么}：十种题型里，前四种在别的节讲得少，是复习课新加的；后六种是 \S2--\S7 的规则换了个考法。做题顺序永远是\textbf{先判「这一题在问什么」→ 再回对应节的规则}，不要凭语感。

\textbf{逐行读法}：第 1、4 两题是「选择题／填空题」，只需要认形式；第 2、3 两题最容易失分，因为四个连词都可能被译成「和／因为」；第 10 题组句考的是词序，先摆主语再摆动词，副词（очень、хорошо́、вчера́）放在动词前面。

\aside{复习课件 p33 的组句题给的是「词语一堆、自己排」，这正是 Ⅲ 部分 15 分最可能的形式。练习时把每题的「主语—动词—宾语—副词」四段先划出来，再排顺序，比直接读顺更稳。}
